\chapter{Proof details and auxiliary derivations}
\label{app:proofs}

\section{Residue projector formula}
Let $Q\ket{\psi_k}=k\ket{\psi_k}$. Then
\[
\frac1M\sum_{\ell=0}^{M-1}
\omega_M^{-r\ell}e^{2\pi i\ell Q/M}\ket{\psi_k}
=
\left(
\frac1M\sum_{\ell=0}^{M-1}
\omega_M^{(k-r)\ell}
\right)\ket{\psi_k}.
\]
The bracket equals one when $k\equiv r\pmod M$ and zero otherwise. This proves \cref{eq:residue-projector} and the orthogonality of the residue packets.

\section{Circuit leakage telescoping}
Let $V_j=PG_jP$. Insert and subtract one factor at a time:
\[
G_m\cdots G_1P-V_m\cdots V_1P
=
\sum_{k=1}^m
G_m\cdots G_{k+1}(G_k-V_k)V_{k-1}\cdots V_1P.
\]
Every $G_j$ has norm one and every $V_j$ has norm at most one. On $P\Hh$,
\[
(G_k-V_k)P=(I-P)G_kP.
\]
Taking norms gives \cref{prop:sum-leakage}.

\section{Hypercube quotient coupling}
Let
\[
\ket k
=
\binom nk^{-1/2}
\sum_{|x|=k}\ket x.
\]
Every weight-$k$ string has $n-k$ neighbors of weight $k+1$, while every weight-$(k+1)$ string has $k+1$ neighbors of weight $k$. Counting the bipartite edges yields
\[
\bra{k+1}A\ket k
=
\frac{\binom nk(n-k)}
{\sqrt{\binom nk\binom n{k+1}}}
=
\sqrt{(k+1)(n-k)},
\]
proving \cref{eq:hypercube-coupling}.

\section{Symmetric-subspace dimension}
A symmetric basis state of $n$ qudits is indexed by occupation numbers
\[
(n_1,\ldots,n_d),
\qquad
n_1+\cdots+n_d=n.
\]
The stars-and-bars count is
\[
\binom{n+d-1}{d-1},
\]
proving \cref{eq:symmetric-qudit-dim}.

\section{Bosonic and fermionic dimensions}
Bosonic occupations are weak compositions of $N$ into $m$ parts, giving $\binom{N+m-1}{N}$. Fermionic occupations are $N$-element subsets of $m$ modes, giving $\binom mN$.

\section{Acceptance-probability stability}
If normalized states $\ket\psi,\ket\phi$ satisfy
\[
\|\ket\psi-\ket\phi\|\le\varepsilon
\]
and $0\le M\le I$, then
\[
|\bra\psi M\ket\psi-\bra\phi M\ket\phi|
\le 2\varepsilon.
\]
Indeed, add and subtract $\bra\phi M\ket\psi$ and use Cauchy--Schwarz. This converts state-vector leakage bounds into measurement-error bounds.

\section{Orbit quotient invariance}
For an automorphism group $G$, the adjacency matrix commutes with the permutation representation. The fixed subspace of the action is spanned by normalized orbit states. Therefore adjacency preserves that span. Polynomial-time simulation additionally requires direct computation of the quotient entries.
